\section{From an input signal to separation of trace moments}
\label{sec:separation}

We prove Theorem~\ref{thm:separation}.  Proposition~\ref{prop:multiscale}
preserves a detectable input-coordinate signal while reducing its decay
rate.  We first choose enough scales to make that rate at most $1/n$.
The resulting gap parameters need not belong to the disk on which trace
moments directly control the signal.  A simultaneous one-variable
continuation will bridge this difference.

Throughout this section $0<p<1$, and $q_*,Q,Y$ are as in
\eqref{eq:aux-parameters}.  Recall from Section~\ref{sec:setup} that
\begin{equation}
  1\le \mathsf X(n,p)\le Y\le 2\mathsf X(n,p),
  \qquad QY=Q+\log n\le C\log n.
  \label{eq:aux-comparisons}
\end{equation}
\subsection{A scale schedule and the surviving signal}

Set
\begin{equation}
  \mu=7.98,\qquad \alpha=0.001,\qquad
  h=\lceil\log_2Y\rceil+H,\qquad D_*=4^h,
  \label{eq:separation-cutoff}
\end{equation}
where the absolute integer $H$ will be fixed below.  In particular,
$D_*\le4^{H+1}Y^2$.  Fix an arbitrary positive moment functional of
order $2D$, where $D\ge D_*$, with first wrong position $d$,
and let $\Delta$ be its
difference from evaluation at the true word, as in Section~\ref{sec:setup}.
All constants chosen below are independent of this system, $n$, and $p$.

For an absolute constant $\eta>0$, define
\begin{equation}
 S_0=4,\qquad
 S_j=\eta q_*Y^{-\alpha}S_{j-1}
               \min\{1,D_*S_{j-1}\}\quad(1\le j\le h).
 \label{eq:scale-schedule}
\end{equation}
The linear regime of this recurrence brings $D_*S_j$ below one; its
quadratic regime then makes $S_j$ small rapidly.  Both regimes will also
be used to control the arguments of the final gap parameters.

Choose an integer $J_*$ with $\alpha J_*\ge3$, and set $H=J_*+3$.
To verify the decrease precisely, write
\[
 \vartheta=\eta q_*Y^{-\alpha},\qquad t_j=D_*S_j.
\]
Then $t_j=\vartheta t_{j-1}\min\{1,t_{j-1}\}$, so, regardless of when
the quadratic regime begins,
\[
 t_{J_*}\le \vartheta^{J_*}t_0
 \le4^{H+2}\eta^{J_*}Y^{2-\alpha J_*}\le1
\]
once $\eta$ is small enough.  Require also
\[
 \eta\le\min\{\eta_0,1/4,e^{-1}\},
\]
where $\eta_0$ is from Proposition~\ref{prop:multiscale}.
Now $\log(1/\vartheta)\ge Q$, and $t_j\le1$ for $j\ge J_*$.  Iterating
$t_j=\vartheta t_{j-1}^2$ in this range gives
\begin{equation}
 t_h\le\vartheta^{\,2^{h-J_*}-1}
       \le \exp\{-Q(2^{h-J_*}-1)\}\le n^{-1},
 \qquad S_h\le n^{-1}.
 \label{eq:terminal-scale}
\end{equation}
The last inequality follows from $2^{h-J_*}\ge8Y$ and
$Q(Y-1)=\log n$.  At every transition we also have
$0<S_j\le\min\{1,\eta_0S_{j-1}\}$, as required by the proposition.

We need an upper bound on the cost of these small scales as well.
Put $A_0=Q+\log Y$.  The recurrence implies, with
$\ell_j=\max\{0,\log(1/S_j)\}$,
\[
 \ell_j\le2\ell_{j-1}+\log(1/\vartheta),\qquad
 \ell_0=0,\qquad \log(1/\vartheta)\le C A_0.
\]
Consequently
\begin{equation}
 \log(1/S_j)\le C2^jA_0\quad(0\le j\le h).
 \label{eq:scale-cost}
\end{equation}

We can now iterate Proposition~\ref{prop:multiscale}.  Begin with the
order-one feature $f_i^{(0)}=X_i$, at decay $s_0=4$ and frequency zero.
The prefix agreement gives $\Delta[X_i]=0$ for $i<d$, whereas
$\abs{\Delta[X_d]}=1$.  Since $\abs{\Delta[X_i]}\le2$ for every $i$,
\begin{equation}
 \abs{\mathcal S_{f^{(0)}}(4,0)}
 \ge1-2\sum_{r\ge1}e^{-4r}
 =1-\frac{2}{e^4-1}>\frac12.
 \label{eq:initial-signal}
\end{equation}
Use error exponents
\begin{equation}
 E_j=C_5A_0\mu^j\quad(0\le j\le h),
 \label{eq:scale-errors}
\end{equation}
where $C_5$ is an absolute constant chosen after $H$ and $\eta$.

Here are all the induction bounds.  At stage $j$ there is a feature
$f^{(j)}$ of order $k_j\le4^j$, with
\begin{equation}
 \begin{gathered}
 \beta S_j\le s_j\le S_j,
 \qquad \abs{\omega_j}\le2.002s_j,
 \qquad \abs{\mathcal S_{f^{(j)}}(s_j,\omega_j)}\ge e^{-E_j},\\
 \log H_{f^{(j)}}\le C4^jA_0.
 \end{gathered}
 \label{eq:scale-induction}
\end{equation}
The initial case follows from \eqref{eq:initial-signal} and
$H_{f^{(0)}}=1$, on increasing $C_5$.  At the transition from $j-1$ to
$j$, the degree requirement is
\[
 4k_{j-1}\le4^j\le D_*\le D.
\]
The feature-height conclusion of the proposition and
\eqref{eq:scale-cost} give
\[
 \log H_{f^{(j)}}
 \le4\log H_{f^{(j-1)}}
       +3\log\bigl(1+1/(\beta S_j)\bigr)
 \le4\log H_{f^{(j-1)}}+C2^jA_0.
\]
Summing this recurrence proves the last bound in
\eqref{eq:scale-induction}.  To check the exponent hypothesis of the
proposition, its bracket at this transition is at most
\[
 C\bigl(1+j4^{j-1}+4^{j-1}A_0+2^jA_0\bigr).
\]
Because $\mu>4$, the sequences
$j(4/\mu)^{j-1}$ and $2^j/\mu^{j-1}$ are bounded.  Thus one fixed
choice of $C_5$ makes this bracket times $C_0$ at most $E_{j-1}$ for
every $j$.  All hypotheses of Proposition~\ref{prop:multiscale} hold,
and its signal bound is $e^{-\mu E_{j-1}}=e^{-E_j}$.

Write $f=f^{(h)}$, $k=k_h$, $s=s_h$, and $\omega=\omega_h$.
We have obtained a signal of magnitude at least $e^{-E_h}$ with
$k\le D_*$ and $s\le1/n$.  The next step records where its parameters
lie in the unit disk; this information will control the cost of
continuing the signal to the trace domain.

\subsection{Damping under multiplication}

Call the parameters introduced between endpoints in each application of
Proposition~\ref{prop:multiscale} \emph{elementary parameters}.  Include
also the final outer parameter
\[
 z_0=e^{-s+\ii\omega}.
\]
An elementary parameter $e^{-a+\ii\theta}$ introduced at stage $j$
satisfies $a\ge\beta S_j$ and $\abs{\theta}\le C_1S_{j-1}$.
Using \eqref{eq:scale-schedule}, choose a sufficiently small absolute
$c_2>0$ so that all these parameters obey
\begin{equation}
 a\ge\gamma\min\{\abs{\theta},D_*\theta^2\},
 \qquad \gamma=c_2q_*Y^{-\alpha}\le1.
 \label{eq:elementary-damping}
\end{equation}
For an explicit verification, the increasing function
$u\mapsto u\min\{1,D_*u\}$ changes by at most a factor
$\max\{C_1,C_1^2\}$ when $u$ is replaced by $C_1u$.
It therefore suffices to take
$c_2\le\beta\eta/\max\{C_1,C_1^2\}$.
Shrinking $c_2$ further covers $z_0$, since
$\abs{\omega}\le2.002s$.  Conjugating any elementary parameter
preserves \eqref{eq:elementary-damping}.

Each new inner gap contains at most one parameter from each of at most
four copies of the preceding feature, and at most one new elementary
parameter.  Count repeated factors with multiplicity; conjugation does
not change the count.  Starting with no inner parameters, the maximum
number $N_j$ of elementary factors in a stage-$j$ inner gap satisfies
\[
 N_0=0,\qquad N_j\le4N_{j-1}+1,
 \qquad N_j\le\frac{4^j-1}{3}.
\]
In particular, an inner gap of $f$, even after multiplication by $z_0$,
contains at most $(4^h+2)/3\le D_*$ elementary factors.

The following elementary lemma explains why this count matters.  The
Cayley map $z\mapsto(1+z)/(1-z)$ sends the unit disk to the right
half-plane; we need a quantitative distance from its boundary.

\begin{lemma}[Products with controlled damping]
\label{lem:product-cayley}
Let $N\ge1$ be an integer and $0<\gamma\le1$.  Suppose $z$ is a
nonempty product of at most $N$ numbers $e^{-a_j+\ii\theta_j}$, where
$a_j>0$, $\theta_j\in\RR$, and
\[
 a_j\ge\gamma\min\{\abs{\theta_j},N\theta_j^2\}.
\]
There is an absolute $c_3>0$ such that
\[
 \re\frac{1+z}{1-z}\ge c_3\gamma.
\]
\end{lemma}

\begin{proof}
Write $z=e^{-T+\ii\psi}$, where $T=\sum_j a_j>0$ and
$\psi=\sum_j\theta_j$.  Let
\[
 U=\sum_{\abs{\theta_j}\ge1/N}\abs{\theta_j},
 \qquad V=\sum_{\abs{\theta_j}<1/N}\abs{\theta_j}.
\]
Summing the hypotheses and applying Cauchy--Schwarz to at most $N$
small angles yields
\[
 T\ge\gamma\left(U+N\sum_{\abs{\theta_j}<1/N}\theta_j^2\right)
   \ge\gamma(U+V^2).
\]
If $U\ge1$ or $V\ge1$, the last bracket is at least one.
Otherwise $U+V^2\ge U^2+V^2\ge(U+V)^2/2$.
Since $\abs{\psi}\le U+V$, in all cases
\[
 T\ge\tfrac12\gamma\min\{1,\psi^2\}.
\]
For $0<T\le1$ it follows that
\[
 1-\abs z^2\ge cT,\qquad
 \abs{1-z}^2\le C\bigl(T^2+\min\{1,\psi^2\}\bigr)
                 \le C'T/\gamma.
\]
Their ratio is $\re((1+z)/(1-z))$.  If $T>1$, that ratio is at least
$(1-e^{-2})/4$, which also suffices because $\gamma\le1$.
\end{proof}

\subsection{The exact trace generating identity}

Let $z_1,\ldots,z_{k-1}$ denote the inner gap parameters of the final
feature $f$.  Define the polynomial in $k$ complex variables
\begin{equation}
 F(Z_0,\ldots,Z_{k-1})
 =\sum_{1\le i_1<\cdots<i_k\le n}
  \Delta\left[\prod_{r=1}^kX_{i_r}\right]
  Z_0^{i_1}\prod_{r=1}^{k-1}Z_r^{i_{r+1}-i_r}.
 \label{eq:input-polynomial}
\end{equation}
At the point
\[
 Z_0^*=z_0,\qquad Z_r^*=z_rz_0\quad(1\le r<k),
\]
the exponents of $z_0$ telescope to $i_k$, so
\begin{equation}
 F(Z^*)=z_0^d\mathcal S_f(s,\omega),
 \qquad \abs{F(Z^*)}\ge e^{-E_h-1}.
 \label{eq:input-polynomial-signal}
\end{equation}
Here $sd\le nS_h\le1$.  Each coordinate of $Z^*$ is a nonempty
product of at most $D_*$ elementary parameters, and hence
Lemma~\ref{lem:product-cayley} gives
\begin{equation}
 \re\frac{1+Z_r^*}{1-Z_r^*}\ge c_3\gamma\quad(0\le r<k).
 \label{eq:input-cayley-margin}
\end{equation}
We shall also need a bound on the coordinates' collective proximity to
the unit circle.  If $z_r=e^{-a_r+\ii\theta_r}$, then
\[
 \frac1{1-\abs{Z_r^*}}
 \le1+\frac1{a_r+s}\le1+\frac1{a_r}\quad(r\ge1),
 \qquad
 \frac1{1-\abs{Z_0^*}}\le1+\frac1s.
\]
These inequalities follow from $e^u\ge1+u$ for $u>0$.
The height bound in \eqref{eq:scale-induction} and
\eqref{eq:scale-cost} therefore imply
\begin{equation}
 W_*:=\log\prod_{r=0}^{k-1}\frac1{1-\abs{Z_r^*}}
 \le\log H_f+\log\bigl(1+1/(\beta S_h)\bigr)
 \le C4^h A_0.
 \label{eq:input-boundary-cost}
\end{equation}

We next express the same polynomial in terms of observable trace
moments on a different domain.  This uses the multivariate trace
generating identity of~\cite[Section~5.2]{CDLSS2021Smoothed}; see also
\cite[Fact~20]{CDLS2024}.  We give the required counting argument
directly.  For a rank tuple
$\ell=(\ell_1,\ldots,\ell_k)$ with
$1\le\ell_1<\cdots<\ell_k\le n$, abbreviate
$T_\ell=T_{\{\ell_1,\ldots,\ell_k\}}$.
For $u=(u_0,\ldots,u_{k-1})\in\CC^k$, set
\[
 G(u)=\sum_{1\le\ell_1<\cdots<\ell_k\le n}
       \Delta[T_\ell]\,
       u_0^{\ell_1-1}
       \prod_{r=1}^{k-1}u_r^{\ell_{r+1}-\ell_r-1}.
\]
Then, with $Z_r=q+pu_r$, the exact polynomial identity is
\begin{equation}
 \left(\prod_{r=0}^{k-1}Z_r\right)G(u)=p^kF(Z).
 \label{eq:trace-generating-identity}
\end{equation}
To prove it, fix input positions $i_1<\cdots<i_k$ that supply the
specified retained bits.  Their own retention contributes $p^k$.
Every input position before $i_1$ contributes $q+pu_0$, according to
whether it is deleted or retained; there are $i_1-1$ such positions.
Similarly, the $i_{r+1}-i_r-1$ positions in the $r$th intervening gap
each contribute $q+pu_r$.  Positions after $i_k$ contribute a total
factor of one.  The retained counts in these intervals are exactly
$\ell_1-1$ and $\ell_{r+1}-\ell_r-1$.  Summing over all input tuples
and applying $\Delta$ gives $G(u)$ as $p^k$ times the polynomial in
\eqref{eq:input-polynomial} with every displayed gap exponent reduced
by one.  Multiplication by $\prod_rZ_r$ proves
\eqref{eq:trace-generating-identity}, including points where some
$Z_r=0$.

Suppose, for the moment, that
\begin{equation}
 \abs{\Delta[T_I]}\le\varepsilon
 \quad\text{for every }I\subseteq\{1,\ldots,n\}
 \text{ with }1\le\abs I\le D_*,
 \label{eq:trace-closeness}
\end{equation}
where $\varepsilon>0$.  Since $k\le D_*$, the rank sum has at most
$n^k$ terms, so $\abs{G(u)}\le n^k\varepsilon$ when all
$\abs{u_r}\le1$.  On this polydisk $\abs{Z_r}\le q+p=1$.
Thus \eqref{eq:trace-generating-identity} gives
\begin{equation}
 \abs{F(Z)}\le (n/p)^k\varepsilon
 \qquad\text{if }\abs{Z_r-q}\le p\quad(0\le r<k).
 \label{eq:trace-polydisk-bound}
\end{equation}
This calculation uses only ordinary trace ranks and the convention that
a rank product is zero when the trace is too short.

\subsection{Simultaneous continuation and the separation bound}

The lower bound \eqref{eq:input-polynomial-signal} and the upper bound
\eqref{eq:trace-polydisk-bound} concern different points.  In Cayley
coordinates the disk $\abs{Z-q}\le p$ becomes a half-plane:
\begin{equation}
 \abs{Z-q}\le p
 \quad\Longleftrightarrow\quad
 \re\frac{1+Z}{1-Z}\ge\frac qp
 \qquad(\abs Z<1).
 \label{eq:cayley-trace-disk}
\end{equation}
Indeed, writing $w=(1+Z)/(1-Z)$ and expanding
$\abs{pw-(1+q)}^2\le p^2\abs{w+1}^2$ reduces it to
$4q\le4p\re w$.  We will move all coordinates to this half-plane
using one complex variable.  A single strip interpolation avoids
incurring a separate continuation loss for every coordinate.

Write
\[
 w_r=\frac{1+Z_r^*}{1-Z_r^*}=\xi_r+\ii\zeta_r,
 \qquad \xi_r\ge c_3\gamma,
\]
where $\xi_r,\zeta_r\in\RR$.  For $v\in\CC$ set
\begin{equation}
 w_r(v)=(1+v)\xi_r+\ii\zeta_r,
 \qquad Z_r(v)=\frac{w_r(v)-1}{w_r(v)+1},
 \qquad
 R=1+\frac{q}{pc_3\gamma}.
 \label{eq:simultaneous-path}
\end{equation}
Figure~\ref{fig:simultaneous-cayley} shows how a common real value of
$v$ moves two possible starting coordinates into the trace half-plane.
The proof uses the analytic maps \eqref{eq:simultaneous-path} for complex
$v$ throughout a strip.

\suppressfloats[t]
\begin{figure}[htbp]
\centering
\begin{tikzpicture}[x=1.2cm,y=1.2cm,>=Stealth,font=\small]
 \begin{scope}
  \fill[blue!4] (0,0) circle (1);
  \fill[teal!17] (0.4,0) circle (0.6);
  \draw[black!35,->] (-1.2,0)--(1.32,0) node[right] {$\re Z$};
  \draw[black!35,->] (0,-1.16)--(0,1.2);
  \draw[black!65] (0,0) circle (1);
  \draw[teal!65!black] (0.4,0) circle (0.6);
  \fill (0.4,0) circle (0.022) node[below right] {$q$};
  \node[below right] at (1,0) {$1$};
  \fill[blue!70!black] (-0.2027,0.5791) circle (0.027)
       node[above left] {$Z_0^*$};
  \fill[orange!80!black] (-0.1753,-0.4310) circle (0.027)
       node[below left] {$Z_1^*$};
  \node at (0,1.5) {The $Z$-plane};
  \node at (0,-1.42) {$\abs{Z-q}\le p$ is shaded};
 \end{scope}
 \draw[->,thick] (1.68,0.22)--(3.78,0.22)
       node[midway,above=5pt] {$w=\dfrac{1+Z}{1-Z}$};
 \begin{scope}[xshift=6cm]
  \fill[blue!4] (0,-1.1) rectangle (2.85,1.1);
  \fill[teal!17] (0.6667,-1.1) rectangle (2.85,1.1);
  \draw[black!35,->] (-0.15,0)--(3,0) node[right] {$\re w$};
  \draw[black!65,->] (0,-1.16)--(0,1.22);
  \draw[dashed,teal!65!black] (0.6667,-1.1)--(0.6667,1.1);
  \node[below right] at (0.6667,0) {$q/p$};
  \draw[->,thick,blue!70!black] (0.35,0.65)--(1.6333,0.65);
  \fill[blue!70!black] (0.35,0.65) circle (0.027)
        node[above left] {$w_0$};
  \fill[blue!70!black] (1.6333,0.65) circle (0.027)
        node[above] {$w_0(R)$};
  \draw[->,thick,orange!80!black] (0.5,-0.55)--(2.3333,-0.55);
  \fill[orange!80!black] (0.5,-0.55) circle (0.027)
        node[below left] {$w_1$};
  \fill[orange!80!black] (2.3333,-0.55) circle (0.027)
        node[below] {$w_1(R)$};
  \node at (1.4,1.5) {The $w$-plane};
  \node at (1.4,-1.42) {Both paths use the same $0\le v\le R$};
 \end{scope}
\end{tikzpicture}
\caption{Schematic geometry of the simultaneous continuation.
The Cayley map turns the deletion disk into
$\re w\ge q/p$.  Two possible starting coordinates are shown.
The horizontal arrows have different lengths but share the same
parameter: $w_r(v)=(1+v)\re w_r+\ii\operatorname{Im}w_r$.
Only their real traces $0\le v\le R$ are drawn.  The proof interpolates
$F(Z_0(v),\ldots,Z_{k-1}(v))$ on a single complex strip, simultaneously
for every coordinate.}
\label{fig:simultaneous-cayley}
\end{figure}

On the closed strip $-1/2\le\re v\le R$ we have
$\re w_r(v)=(1+\re v)\xi_r>0$, so all $Z_r(v)$ lie in the unit
disk and depend analytically on $v$.  At $v=0$ they equal $Z_r^*$.
On the right boundary,
\[
 \re w_r(v)=(1+R)\xi_r\ge q/p,
\]
and therefore \eqref{eq:trace-polydisk-bound} applies.  Moreover
\begin{equation}
 1+R\le C\frac{Y^\alpha}{p},
 \label{eq:strip-width}
\end{equation}
because $\gamma=c_2q_*Y^{-\alpha}$ and $q\le q_*$.

We need a bound on the left boundary that does not involve $n^k$.
The identity $1-\abs{Z}^2=4\re w/\abs{w+1}^2$ gives, when
$\re v=-1/2$,
\begin{equation}
 1-\abs{Z_r(v)}
 \ge\frac{1-\abs{Z_r^*}}{4(1+\abs v)^2}.
 \label{eq:left-boundary-distance}
\end{equation}
To check the constant, first use
$1-\abs{Z_r(v)}\ge \xi_r/\abs{w_r(v)+1}^2$ and
$1-\abs{Z_r^*}\le4\xi_r/\abs{w_r+1}^2$.
Then use
\[
 \abs{w_r(v)+1}\le\abs{w_r+1}+\xi_r\abs v
                   \le(1+\abs v)\abs{w_r+1}.
\]
The coefficients of $F$ have magnitude at most two.  Summing
independent geometric series over the positive prefix and gap lengths
therefore gives, for every point in the open unit polydisk,
\[
 \abs{F(Z)}\le2\prod_{r=0}^{k-1}\frac1{1-\abs{Z_r}}.
\]
It follows from \eqref{eq:left-boundary-distance} that
\[
 g(v):=\frac{F(Z_0(v),\ldots,Z_{k-1}(v))}{(2+v)^{2k}}
\]
has the boundary bounds
\begin{equation}
 \begin{aligned}
 \log\abs{g(v)}&\le W_*+Ck+C
                    &&(\re v=-1/2),\\
 \log\abs{g(v)}&\le k\log(n/p)+\log\varepsilon
                    &&(\re v=R).
 \end{aligned}
 \label{eq:strip-side-bounds}
\end{equation}
On the left, $(1+\abs v)/\abs{2+v}$ is bounded by an absolute
constant.  On the right, $\abs{2+v}\ge2+R>1$.
The function $g$ is analytic and bounded on the strip: $F$ is a finite
polynomial evaluated in the unit polydisk, and $\abs{2+v}\ge3/2$.

For clarity, the strip estimate we use is the following elementary
consequence of the maximum principle.  If a bounded analytic function
on $a\le\re v\le b$ has logarithmic boundary bounds $L_-,L_+$,
then at a real $t\in[a,b]$ its logarithm is at most
$((b-t)L_-+(t-a)L_+)/(b-a)$.  Multiply the function by the exponential
of the negative affine interpolant to make both side bounds one.
On rectangles of height tending to infinity, apply the maximum
principle after multiplying further by
$\exp\{\tau(v^2-\max\{a^2,b^2\})\}$, for $\tau>0$.
This factor has modulus at most one on the vertical sides and tends
uniformly to zero on the horizontal sides.  Let the height tend to
infinity and then $\tau$ tend to zero.  This proves the estimate,
including zeros of the function by continuity.

At $t=0$ in our strip, the right boundary has weight $1/(2R+1)$.
By \eqref{eq:input-polynomial-signal},
$\log\abs{g(0)}\ge-E_h-1-2k\log2$.
Substituting \eqref{eq:strip-side-bounds} consequently gives
\begin{align}
 \log(1/\varepsilon)
 &\le k\log(n/p)+(2R+1)(E_h+1+2k\log2)
             +2R(W_*+Ck+C)\notag\\
 &\le k\log(n/p)
       +C(1+R)\bigl(E_h+4^hA_0\bigr).
 \label{eq:separation-exponent}
\end{align}
This is the desired quantitative transfer from input coordinates to
trace coordinates.

It remains to compare the exponent with the budget in
\eqref{eq:budget}.  Put $t=\log_2\mu$.  The choices in
\eqref{eq:separation-cutoff} satisfy $t+\alpha<3$, and
\[
 4^h\le CY^2,\qquad E_h\le CA_0Y^t,
 \qquad A_0\le Q(1+\log Y),\qquad QY\le C\log n.
\]
Thus \eqref{eq:strip-width} yields
\begin{align*}
 (1+R)E_h
 &\le \frac Cp(\log n)Y^{t+\alpha-1}(1+\log Y)
 \le \frac Cp(\log n)Y^2(1+\log Y),\\
 (1+R)4^hA_0
 &\le \frac Cp(\log n)Y^{1+\alpha}(1+\log Y)
 \le \frac Cp(\log n)Y^2(1+\log Y).
\end{align*}
Also, since $n\ge2$ and $\log(1/p)\le1/p$,
\[
 k\log(n/p)\le CY^2\bigl(\log n+\log(1/p)\bigr)
                \le \frac Cp(\log n)Y^2.
\]
Using \eqref{eq:aux-comparisons}, we conclude that the right-hand side
of \eqref{eq:separation-exponent} is at most
\begin{equation}
 \frac Cp(\log n)Y^2(1+\log Y)
 \le C_6 p^{-1}(\log n)\mathsf X(n,p)^2
                  (1+\log\mathsf X(n,p))^6
 =C_6B(n,p).
 \label{eq:final-B-bound}
\end{equation}
Taking $\varepsilon=e^{-C_{\mathrm{sep}}B(n,p)}$ with
$C_{\mathrm{sep}}>C_6$ contradicts
\eqref{eq:separation-exponent}.  Hence
\eqref{eq:trace-closeness} fails for that $\varepsilon$, proving
Theorem~\ref{thm:separation} with the cutoff
\eqref{eq:separation-cutoff}.

All constants were fixed in the following order: those of
Proposition~\ref{prop:multiscale}; then $J_*,H$ and $\eta$; then
$C_5,c_2$; and finally $C_{\mathrm{sep}}$.  There is no dependence on
the denominator of $p$ or on how close it is to an endpoint.
In particular, when $q<1/n$ the cutoff $q_*=1/n$ keeps $Y<2$ and
$Q=1+\log n$, while $q/q_*\le1$ keeps the strip estimate uniform.
When $p$ is small, the strip width is controlled by the single factor
$p^{-1}$ already present in the budget.  The case $p=1$ will be handled
directly in the decoding argument.
